\documentclass[border=10pt]{standalone}
\usepackage[no-math]{fontspec}
\usepackage[brazil]{babel}
\usepackage{amsmath,amssymb}
\IfFileExists{C:/Windows/Fonts/Lato_Regular.ttf}{%
  \setmainfont{Lato}[
    Path = C:/Windows/Fonts/,
    UprightFont = Lato_Regular.ttf,
    BoldFont = Lato_Bold.ttf,
    ItalicFont = Lato_Italic.ttf,
    BoldItalicFont = Lato_BoldItalic.ttf
  ]
  \setsansfont{Lato}[
    Path = C:/Windows/Fonts/,
    UprightFont = Lato_Regular.ttf,
    BoldFont = Lato_Bold.ttf,
    ItalicFont = Lato_Italic.ttf,
    BoldItalicFont = Lato_BoldItalic.ttf
  ]
}{%
  \RequirePackage[default]{lato}
}
\renewcommand{\familydefault}{\sfdefault}

\usepackage{xcolor}
\usepackage{tikz}
\usepackage{pgfplots}
\pgfplotsset{compat=1.18}
\usetikzlibrary{arrows.meta,positioning,shapes.geometric,fit,backgrounds,calc,decorations.pathreplacing,patterns}

% Paleta institucional UEFS / Sociologia Rural
\definecolor{srNavy}{HTML}{1F4268}
\definecolor{srNavyDark}{HTML}{152E4A}
\definecolor{srGold}{HTML}{B38A3D}
\definecolor{srGoldDark}{HTML}{8C6927}
\definecolor{srTeal}{HTML}{0E7C7B}
\definecolor{srForest}{HTML}{2D6A4F}
\definecolor{srAlert}{HTML}{9E2A2B}
\definecolor{srInk}{HTML}{26313D}
\definecolor{srMuted}{HTML}{5B6573}

\definecolor{bgNavy}{HTML}{EDF2F7}
\definecolor{bgGold}{HTML}{F8F3E9}
\definecolor{bgTeal}{HTML}{EBF6F5}
\definecolor{bgForest}{HTML}{EEF5F0}
\definecolor{bgAlert}{HTML}{FDF2F2}

\tikzset{
  >=Stealth,
  srBox/.style={
    draw=#1, fill=white, rounded corners=5pt, line width=1.15pt,
    inner sep=7pt, align=center, font=\footnotesize\sffamily, text=srInk
  },
  srTint/.style={
    draw=#1, fill=#1!8!white, rounded corners=5pt, line width=1.15pt,
    inner sep=7pt, align=center, font=\footnotesize\sffamily, text=srInk
  },
  srBadge/.style={
    fill=#1, text=white, rounded corners=9pt, inner xsep=8pt, inner ysep=3pt,
    font=\scriptsize\bfseries\sffamily
  },
  srTitle/.style={font=\normalsize\bfseries\sffamily, text=srNavy, align=center},
  srSub/.style={font=\footnotesize\sffamily, text=srMuted, align=center},
  srArrow/.style={->, >=Stealth, line width=1.35pt},
  srDArrow/.style={<->, >=Stealth, line width=1.35pt}
}


\begin{document}
\begin{tikzpicture}
  \begin{axis}[
    xbar,
    width=11.4cm,
    height=6.2cm,
    at={(0cm,0cm)},
    xmin=0,
    xmax=112,
    ymin=-0.55,
    ymax=5.55,
    bar width=16pt,
    bar shift=0pt,
    axis lines*=left,
    axis line style={draw=srMuted!65, line width=0.8pt},
    major tick style={draw=none},
    xmajorgrids=true,
    grid style={dashed, draw=srMuted!35, line width=0.5pt},
    xtick={0,20,40,60,80,100},
    xticklabels={0\%,20\%,40\%,60\%,80\%,100\%},
    xticklabel style={font=\footnotesize\sffamily, text=srMuted},
    ytick={0,1,2,3,4,5},
    yticklabels={Internet fixa, Rádio, YouTube, WhatsApp, Celular, TV Aberta},
    yticklabel style={font=\small\bfseries\sffamily, text=srInk, xshift=-3pt},
    nodes near coords={\pgfmathprintnumber[fixed,precision=0,assume math mode=true]{\pgfplotspointmeta}\%},
    nodes near coords align={horizontal},
    every node near coord/.append style={
      font=\small\bfseries\sffamily,
      text=srInk,
      fill=white,
      fill opacity=0.88,
      text opacity=1,
      inner sep=1.5pt,
      xshift=4pt
    },
    clip=false
  ]
    % Barras em srNavy (infraestrutura fixa e radiodifusão sonora)
    \addplot[fill=srNavy, draw=srNavyDark, line width=0.6pt] coordinates {
      (21, 0)
      (43, 1)
    };
    % Barras em srTeal (plataformas digitais)
    \addplot[fill=srTeal, draw=srTeal!80!black, line width=0.6pt] coordinates {
      (45, 2)
      (79, 3)
    };
    % Barras em srForest (meios de maior capilaridade)
    \addplot[fill=srForest, draw=srForest!80!black, line width=0.6pt] coordinates {
      (87, 4)
      (95, 5)
    };
  \end{axis}

  % Título e subtítulo posicionados logo acima do eixo
  \node[srTitle, font=\large\bfseries\sffamily] (title) at (4.4, 5.65)
    {Meios de Comunicação no Meio Rural Brasileiro (PNAD TIC/IBGE, 2022)};
  \node[srSub] (sub) at (4.4, 5.15)
    {Utilização pela população rural (\%) por meio/plataforma};
\end{tikzpicture}
\end{document}
